\chapter{Background}
\label{chap:background}

This chapter provides the theoretical and methodological background for the remainder of the thesis. It begins by introducing the physics of X-ray CT acquisition and the analytical reconstruction pipeline used throughout this work (\Cref{sec:bg_ct_fundamentals}). It then describes the two main sources of image degradation that the proposed corruption protocol aims to reproduce: photon-counting noise arising at low dose (\Cref{sec:bg_ldct_degradation}) and the deterministic physical artifacts that commonly coexist with it in clinical scans (\Cref{sec:bg_artifacts}). Finally, the chapter reviews the deep learning methods that have come to dominate CT image restoration and situates the present thesis within that body of literature (\Cref{sec:bg_dl_restoration,sec:bg_related}).

\section{CT Imaging Fundamentals}
\label{sec:bg_ct_fundamentals}

\subsection{X-ray Attenuation and the Beer--Lambert Law}
\label{subsec:bg_beer_lambert}

In transmission CT, a narrow X-ray beam is passed through the patient and the transmitted intensity is measured by a detector on the opposite side~\cite{Hounsfield1973, KakSlaney1988}. As the beam traverses tissue, photons are removed from the primary beam through photoelectric absorption and Compton scattering. Under the monoenergetic approximation, the measured intensity along a ray path $L$ is described by the Beer--Lambert law,
\begin{equation}
    I(L) \;=\; I_{0}\,\exp\!\left(-\int_{L} \mu(\boldsymbol{r})\,\mathrm{d}\ell\right),
    \label{eq:beer_lambert}
\end{equation}
where $I_0$ is the incident photon intensity, $\mu(\boldsymbol{r})$ is the spatially varying linear attenuation coefficient at position $\boldsymbol{r} \in \mathbb{R}^{2}$,
and the integral is taken along the geometric ray~$L$. Taking the negative logarithm of the normalized intensity yields the line-integral measurement,
\begin{equation}
    p(L) \;=\; -\ln\!\left(\frac{I(L)}{I_{0}}\right) \;=\; \int_{L} \mu(\boldsymbol{r})\,\mathrm{d}\ell,
    \label{eq:line_integral}
\end{equation}
which is the quantity that enters all subsequent reconstruction~\cite{KakSlaney1988, Buzug2008}. For display and clinical interpretation the reconstructed $\mu$~values are typically
mapped to \glspl{hu} relative to water, $\mathrm{HU}(\boldsymbol{r}) = 1000\,\bigl(\mu(\boldsymbol{r}) - \mu_{\text{water}}\bigr)/\mu_{\text{water}}$~\cite{Hounsfield1973}.

\subsection{The Radon Transform and Sinogram Formation}
\label{subsec:bg_radon}

A complete tomographic acquisition consists of line-integral measurements for a dense set of rays, parameterized by their orientation~$\theta \in [0,\pi)$ and signed
perpendicular offset~$t \in \mathbb{R}$. The mapping from the attenuation image $\boldsymbol{x}(\boldsymbol{r}) = \mu(\boldsymbol{r})$ to the collection of all such line
integrals is the two-dimensional Radon transform~\cite{Radon1917, KakSlaney1988},
\begin{equation}
    \boldsymbol{s}(\theta, t) \;=\; (\mathcal{R}\boldsymbol{x})(\theta, t)
    \;=\; \int_{-\infty}^{\infty}\!\!\int_{-\infty}^{\infty} \boldsymbol{x}(\boldsymbol{r})\,
        \delta\bigl(t - \boldsymbol{r}\!\cdot\!\boldsymbol{n}_{\theta}\bigr)\,\mathrm{d}\boldsymbol{r},
    \label{eq:radon}
\end{equation}
with $\boldsymbol{n}_{\theta} = (\cos\theta, \sin\theta)^{\!\top}$. The discretized representation of $\boldsymbol{s}$, arranged as a two-dimensional array indexed by projection angle and detector position, is referred to as the \emph{sinogram}. This name reflects the fact that a point in image space traces a sinusoidal curve in the projection domain.

CT systems are commonly modeled using either parallel-beam geometry, in which all rays acquired at a given projection angle are mutually parallel, or fan-/cone-beam geometry, in which rays diverge from a point source. In this thesis, we follow the LoDoPaB-CT benchmark and assume a parallel-beam acquisition model throughout~\cite{Leuschner2021}. Specifically, LoDoPaB-CT provides $1000$uniformly spaced projections over $[0,\pi)$, each sampled by 513 detector bins, and reconstructs images of size $362\times362$ at a suitable resolution in the plane. Adopting this fixed geometry simplifies the corruption framework developed in \Cref{chap:corruption}, as each artifact model can be formulated directly in sinogram coordinates $(\theta, t)$.

\subsection{Filtered Back-Projection}
\label{subsec:bg_fbp}

The Fourier slice theorem connects the Radon transform to the 2D Fourier transform of the image: the 1D Fourier transform of a single projection $\boldsymbol{s}(\theta,\cdot)$ in the detector variable equals a radial slice of $\hat{\boldsymbol{x}}(\boldsymbol{\xi})$ at angle $\theta$~\cite{KakSlaney1988}. Inverting this relationship in
polar coordinates gives the \gls{fbp} reconstruction formula,
\begin{equation}
    (\mathcal{R}^{-1}\boldsymbol{s})(\boldsymbol{r}) \;=\; \int_{0}^{\pi}\!
        \bigl(\boldsymbol{s}(\theta, \cdot) * h\bigr)\!\bigl(\boldsymbol{r}\!\cdot\!\boldsymbol{n}_{\theta}\bigr)\,\mathrm{d}\theta,
    \label{eq:fbp}
\end{equation}
where $h$ is a 1D ramp filter (Ram--Lak) whose frequency response is $|\xi|$, optionally apodised by a Hann or Shepp--Logan window to suppress high-frequency noise. Equation~\eqref{eq:fbp} states that an image is recovered by filtering each projection with $h$ and back-projecting (smearing) it across the image plane along the corresponding ray direction. \gls{fbp} is non-iterative, computationally efficient, and remains the analytical reconstruction algorithm in routine clinical use~\cite{KakSlaney1988, Buzug2008}; it is adopted as the fixed reconstruction step throughout this thesis, both for the baseline and as the input pipeline to all learned models.

\section{Low-Dose CT Degradation}
\label{sec:bg_ldct_degradation}

\subsection{Poisson Statistics of Photon Counting}
\label{subsec:bg_poisson}

X-ray photon emission and detection are independent stochastic events, so the number of photons recorded per detector element in a fixed integration window follows
Poisson statistics~\cite{Leuschner2021, Sadia2024}. Combining this with the Beer--Lambert law of \Cref{eq:beer_lambert}, the measured photon count along ray~$L$ is
\begin{equation}
    N(L) \;\sim\; \mathrm{Poisson}\!\Bigl(I_{0}\exp\!\bigl(-\boldsymbol{s}(L)\bigr)\Bigr),
    \label{eq:poisson}
\end{equation}
with $\boldsymbol{s}(L) = \int_{L}\mu\,\mathrm{d}\ell$ the noiseless line integral and $I_0$ the mean number of incident photons per ray. Since the variance of a
Poisson variable equals its mean, the noise level scales as $\sqrt{I_0\,e^{-\boldsymbol{s}}}$ in absolute terms, but the \emph{relative} noise, which propagates into the negative-log sinogram scales as $1/\sqrt{I_0\,e^{\boldsymbol{s}}}$. Two consequences follow: noise is signal-dependent (rays through dense tissue are noisier), and reducing $I_0$ to lower the patient dose increases relative noise everywhere by the inverse square root of the dose-reduction factor. The LoDoPaB benchmark adopts $I_0 = 4096$ photons per ray to emulate a clinically meaningful low-dose regime~\cite{Leuschner2021}; this value is also used throughout the present work.

A useful approximation for analysis and simulation is to treat the negative-log measurement as Gaussian with signal-dependent variance,
\begin{equation}
    \tilde{\boldsymbol{s}}(L) \;\approx\; \boldsymbol{s}(L) \;+\; \varepsilon(L),
    \qquad \varepsilon(L)\,|\,\boldsymbol{s}(L) \;\sim\; \mathcal{N}\!\Bigl(0,\;\tfrac{1}{I_{0}}\,\exp\!\bigl(\boldsymbol{s}(L)\bigr)\Bigr),
    \label{eq:gaussian_approx}
\end{equation}
which is accurate when $I_{0}\,e^{-\boldsymbol{s}}$ is large and is widely used to derive statistically weighted iterative
reconstructions~\cite{Whiting2006, Sadia2024}. We use the exact Poisson model of \Cref{eq:poisson} for simulation in \Cref{chap:corruption}.

\subsection{Noise Propagation through FBP}
\label{subsec:bg_noise_propagation}

Because \gls{fbp} is a linear operator, sinogram noise propagates additively to the reconstructed image. However, the ramp filter amplifies high spatial frequencies and therefore the high-frequency components of the noise as well, so that white noise in the sinogram becomes spatially correlated streak and granular-pattern noise after
back-projection~\cite{KakSlaney1988, Leuschner2021}. This colored noise is qualitatively different from the additive white Gaussian noise that classical denoisers such
as \gls{bm3d}~\cite{Dabov2007} are designed for, and it is one of the reasons why specialized CT denoisers~\cite{Chen2017, Komolafe2024} consistently outperform
generic image-domain denoisers on \gls{ldct} data.

\section{Physical CT Artifacts}
\label{sec:bg_artifacts}

In addition to photon-counting noise, real acquisitions are degraded by deterministic physical effects whose statistics are entirely different from those of
\Cref{eq:poisson}~\cite{Boas2012}. The corruption protocol of \Cref{chap:corruption} models three of the most prevalent clinically: motion, ring, and metal artifacts.

\subsection{Motion Artifacts}
\label{subsec:bg_motion}

A CT acquisition spans on the order of $0.5$ to several seconds; any anatomical motion within that window violates the implicit assumption that the projections at different angles correspond to a single static volume. Respiratory motion, cardiac motion, and involuntary patient movement all introduce projection-time inconsistencies that, after \gls{fbp}, manifest as blurring, doubling, and ghosting in the neighborhood of moving structures~\cite{LuMackie2002, Lu2005, Boas2012}. From a sinogram-domain perspective, motion can be modeled as a time-dependent geometric warp applied to a subset of projections~\cite{Mao2007}, which is the formulation adopted in \Cref{chap:corruption}.

\subsection{Ring Artifacts}
\label{subsec:bg_ring}

Ring artifacts arise from systematic miscalibration or partial failure of individual detector elements~\cite{Sijbers2004, Prell2009}. A constant or slowly varying offset along a single sinogram column (i.e.\ a fixed detector channel) corresponds, after back-projection, to a circle in the image plane with radius equal to the geometric distance from the rotation axis to that detector channel. The result is a pattern of concentric rings centered on the rotation axis that is particularly damaging in the central field of view, where small lesions can be obscured. Classical correction methods include wavelet--Fourier stripe filtering of the sinogram~\cite{Munch2009}. We model these in \Cref{chap:corruption} as additive per-channel gain perturbations.

\subsection{Metal Artifacts}
\label{subsec:bg_metal}

Metallic implants such as hip prostheses, dental fillings, surgical clips, and pacemaker leads have linear attenuation coefficients orders of magnitude greater than soft tissue. Two coupled mechanisms then conspire to produce severe streaks~\cite{ZhangYu2018, Gjesteby2016}. First, photon starvation: along the projection rays that traverse the
metal, $I_0\,e^{-\boldsymbol{s}}$ becomes very small, the Poisson noise in \Cref{eq:poisson} dominates the signal, and the negative-log measurement becomes effectively random. Second, beam hardening: in a polychromatic beam, low-energy photons are preferentially absorbed by the metal, and the residual high-energy spectrum violates the monoenergetic assumption underlying \Cref{eq:beer_lambert}, biasing the line-integral measurement. After \gls{fbp}, the affected rays produce alternating bright and dark streaks radiating from the implant. Decades of work on \gls{mar} have proposed sinogram inpainting, normalized reprojection (NMAR)~\cite{Meyer2010NMAR}, and dual-energy techniques~\cite{Gjesteby2016}, but the artifact remains an open clinical problem.

\section{Deep Learning for CT Image Restoration}
\label{sec:bg_dl_restoration}

\subsection{Image-to-Image Restoration as Supervised Learning}
\label{subsec:bg_supervised}

Modern CT restoration is dominated by supervised deep learning~\cite{LeCun2015, Sadia2024}. Given a paired dataset $\{(\boldsymbol{y}_{i},\boldsymbol{x}_{i})\}_{i=1}^{N}$ of degraded inputs $\boldsymbol{y}_{i}$ (typically the \gls{fbp} of a corrupted sinogram) and reference images $\boldsymbol{x}_{i}$, a network $f(\,\cdot\,;\boldsymbol{\theta})$ is trained to minimize an empirical reconstruction loss,
\begin{equation}
    \boldsymbol{\theta}^{\star}
    \;=\; \arg\min_{\boldsymbol{\theta}} \;
    \frac{1}{N}\sum_{i=1}^{N}
    \mathcal{L}\!\bigl(f(\boldsymbol{y}_{i};\boldsymbol{\theta}),\,\boldsymbol{x}_{i}\bigr),
    \label{eq:supervised_loss}
\end{equation}
where $\mathcal{L}$ is most commonly the $\ell^{1}$ (\gls{mae}) or $\ell^{2}$ (\gls{mse}) pixel loss. We use $\ell^{1}$ throughout, in line with previous CT restoration work that has reported sharper outputs and more stable optimization than $\ell^{2}$~\cite{Chen2017, Zhang2017}.

\subsection{Convolutional Encoder--Decoder Networks}
\label{subsec:bg_unet}

The U-Net architecture~\cite{Ronneberger2015} has become the dominant backbone for image-to-image restoration in medical imaging. It comprises a contracting path, in which strided convolutions and pooling progressively reduce spatial resolution while increasing feature abstraction, a bottleneck at the coarsest scale, and a symmetric expanding path that restores the original image resolution through upsampling. Skip connections link encoder and decoder features at corresponding resolutions, allowing the decoder to recover fine spatial detail that might otherwise be lost during compression. This design makes U-Net particularly well suited to tasks that require pixel-aligned outputs, such as CT denoising and artifact reduction~\cite{Komolafe2024, Zhang2023}. The same encoder--decoder paradigm has also been extended to the output of analytic reconstruction algorithms in sparse-view CT and related inverse problems, with FBPConvNet of Jin et al.~\cite{Jin2017} representing a canonical example.

\subsection{Residual Learning}
\label{subsec:bg_residual}

Two complementary ideas from general deep learning carry over to CT restoration. The first is residual learning at the architectural level: stacking ``identity shortcut'' connections that allow gradients to bypass blocks of layers, originally introduced in ResNet~\cite{He2016}. The second is residual learning at the prediction level:
parameterizing the network to predict the noise residual $\boldsymbol{r} = \boldsymbol{y} - \boldsymbol{x}$ rather than the clean image directly. \gls{dncnn}~\cite{Zhang2017} demonstrated that residual prediction, combined with batch normalization~\cite{Ioffe2015}, substantially accelerates convergence and improves denoising quality, because
the residual is typically much sparser and lower-magnitude than the clean image. \gls{redcnn}~\cite{Chen2017} fused these ideas with the encoder--decoder principle to produce a strong CT-specific baseline that we adopt as a comparison method in \Cref{chap:experiments}. Adversarial training with a Wasserstein objective and a perceptual loss term~\cite{Yang2018WGAN} has since been used to address the over smoothing typical of pixel-only objectives, while model-based unrolled networks
(e.g.\ MoDL~\cite{Aggarwal2019MoDL}, the variational network of Hammernik et al.~\cite{Hammernik2018VarNet}) embed the forward operator directly into the architecture.

\subsection{Cascade and Multi-Stage Approaches}
\label{subsec:bg_cascades}

A natural extension of single-network restoration is to apply several networks in sequence, each refining the output of the previous one. The idea was popularized in medical imaging by the deep cascade of CNNs for dynamic MR reconstruction~\cite{Schlemper2018} and has since been adopted in CT under various formulations~\cite{Wang2020, Zhang2022}. Formally, a $K$-stage cascade is defined by
\begin{equation}
    \boldsymbol{x}_{0} = \boldsymbol{y}, \qquad
    \boldsymbol{x}_{k} \;=\; f_{k}\!\bigl(\boldsymbol{x}_{k-1};\boldsymbol{\theta}_{k}\bigr),
    \quad k = 1,\ldots,K,
    \label{eq:cascade_generic}
\end{equation}
with the final output $\hat{\boldsymbol{x}} = \boldsymbol{x}_{K}$. The choice of how to optimize the parameters $\{\boldsymbol{\theta}_{k}\}$ defines the cascade variant.
Three formulations are studied in \Cref{chap:architectures} and previewed here:
\begin{itemize}
    \item \emph{Independent stage cascade.} All stage parameters $\{\boldsymbol{\theta}_{k}\}$ are optimized simultaneously, but a stop-gradient operator is applied to $\boldsymbol{x}_{k-1}$ so that the loss of stage $k$ does not propagate back to stage $k{-}1$. Each stage is therefore updated only by its own loss, which separates the optimization problems while still allowing all stages to train in parallel within the same pass.
    
    \item \emph{End-to-end cascade.} As in the independent-stage formulation, each stage is associated with its own loss function and all stages are trained simultaneously; however, the stop-gradient constraint is removed so that gradients propagate freely throughout the cascade. This allows the parameters of Stage~1 to adapt in response to the final Stage~2 objective. To reduce the risk of weak gradient signals on the longer optimization path, the per-stage losses are weighted equally in a deep-supervision scheme~\cite{He2016, LeCun2015}.
    
    \item \emph{Residual cascade.} Each stage predicts an additive correction to its input such that $\boldsymbol{x}_{k} = \boldsymbol{x}_{k-1} + f_{k}(\boldsymbol{x}_{k-1}; \boldsymbol{\theta}_{k})$. This formulation inherits the favorable optimization properties of residual learning~\cite{He2016, Zhang2017} while retaining the modular structure of the cascade.
\end{itemize}
\subsection{Dual-Domain Methods}
\label{subsec:bg_dual_domain}

A complementary line of work combines a network operating in the sinogram domain, acting on $\tilde{\boldsymbol{s}}$ prior to \gls{fbp}, with a second network operating in the image domain~\cite{Wang2020, Zhang2022, Lin2019DuDoNet}. The motivation for such dual-domain methods is that the two domains provide complementary information: the sinogram retains measurement-domain structure and data-consistency constraints tied directly to the acquisition process, whereas the image domain is the natural space in which anatomical priors and perceptual structure are expressed. By exploiting both representations jointly, dual-domain approaches can in principle overcome limitations imposed by image-domain restoration alone. Although no dual-domain model is implemented in this thesis, the architectural ceiling identified in \Cref{chap:experiments} provides a quantitative motivation to pursue this direction.

\subsection{Diffusion and Score-Based Generative Models}
\label{subsec:bg_diffusion}

A separate line of research formulates CT restoration as posterior sampling under a learned image prior rather than as direct discriminative regression; here we follow the taxonomy of the survey by Yang et al.~\cite{Yang2024DiffusionSurvey}. Denoising diffusion probabilistic models~\cite{Ho2020} and the closely related score-based formulation~\cite{SongErmon2019, Song2021} learn a generative prior over images through a network \(s_{\boldsymbol{\theta}}(\boldsymbol{x}_t, t)\) trained on progressively perturbed versions of the data. At inference, samples are generated by reversing the diffusion process under the learned score, optionally with additional guidance from a likelihood term that enforces consistency with an observed measurement.

For an inverse problem of the form \(\boldsymbol{y} = \mathcal{A}\boldsymbol{x} + \boldsymbol{n}\), where the forward operator \(\mathcal{A}\) is known, posterior samples from \(p(\boldsymbol{x}\,|\,\boldsymbol{y})\) can be obtained by augmenting the reverse process with a measurement-consistency term~\cite{Song2022, Chung2023}. In CT, \(\mathcal{A}\) corresponds to the discretized Radon transform of \Cref{eq:radon} together with the photon-counting model of \Cref{eq:poisson}. Diffusion posterior sampling (DPS)~\cite{Chung2023} is among the most widely used inference strategies in this setting and is adopted, for example, in the DM4CT framework~\cite{DM4CT2025}, which provides diffusion-based reconstruction baselines for sparse-view, limited-angle, and low-dose CT within a unified codebase. CT-specific examples include score-based reconstruction from sparse measurements~\cite{Song2022} and the limited-angle DOLCE framework~\cite{Liu2023DOLCE}.

Diffusion-based reconstruction offers two attractive features for CT. First, it yields a posterior distribution rather than a single point estimate, which enables uncertainty quantification. Second, because the image prior is learned separately from the measurement model, the same prior can, in principle, be reused across sparse-view, limited-angle, and low-dose settings without retraining. Its principal drawback is inference cost: current samplers typically require tens to hundreds, and sometimes up to thousands, of network evaluations per slice~\cite{Ho2020, Yang2024DiffusionSurvey}, making evaluation on the full LoDoPaB test set substantially more expensive than for the discriminative baselines studied in \Cref{chap:experiments}. In this thesis, diffusion-based methods are therefore treated as an orthogonal family to the cascade study and revisited as a direction for future work in \Cref{chap:conclusion}.

\section{Image Quality Metrics}
\label{sec:bg_metrics}

Two metrics are used throughout this thesis. The first, \gls{psnr}, expresses the per-pixel reconstruction error on a logarithmic scale,
\begin{equation}
    \mathrm{PSNR}(\hat{\boldsymbol{x}},\boldsymbol{x})
    \;=\; 10 \log_{10}\!\frac{\mathrm{MAX}^{2}}{\frac{1}{M}\sum_{j=1}^{M}\bigl(\hat{x}_{j} - x_{j}\bigr)^{2}},
    \label{eq:psnr}
\end{equation}
where $M$ is the number of image pixels and $\mathrm{MAX}$ is the dynamic range used for normalization (the LoDoPaB display window
$[0,\,0.0457]\,\mathrm{mm}^{-1}$~\cite{Leuschner2021}). The second, \gls{ssim}~\cite{Wang2004}, compares local luminance, contrast, and structure between two images and is more strongly correlated with human perception of image quality than \gls{psnr},
\begin{equation}
    \mathrm{SSIM}(\hat{\boldsymbol{x}},\boldsymbol{x})
    \;=\; \frac{(2\mu_{\hat{x}}\mu_{x} + c_{1})(2\sigma_{\hat{x}x} + c_{2})}
              {(\mu_{\hat{x}}^{2} + \mu_{x}^{2} + c_{1})(\sigma_{\hat{x}}^{2} + \sigma_{x}^{2} + c_{2})},
    \label{eq:ssim}
\end{equation}
with $\mu$, $\sigma^{2}$, and $\sigma_{\hat{x}x}$ the local mean, variance, and covariance estimated in a sliding window, and $c_{1}, c_{2}$ small stabilizing constants.
We follow the LoDoPaB-CT convention of computing both metrics per slice and averaging across the test set~\cite{Leuschner2021, Kiss2024}.  Both \gls{psnr} and \gls{ssim}
are known to correlate imperfectly with radiological reading quality, and learned perceptual metrics such as \gls{lpips}~\cite{Zhang2018LPIPS} have been proposed as alternatives; we report \gls{psnr} and \gls{ssim} only, in line with the LoDoPaB benchmarking convention.

\section{Related Work and Positioning}
\label{sec:bg_related}

\Cref{tab:related_work} summarizes the methods most directly relevant to this thesis along three high-level axes: the corruption regime they address, the architectural family to which they belong, and the domain in which they operate. Two broad patterns emerge. First, most learned CT restoration methods are developed for a single isolated degradation type, most commonly Poisson noise at a fixed dose level~\cite{Chen2017, Zhang2017, Komolafe2024, Zhang2023}. Structural artifacts such as metal corruption are typically treated by separate, dedicated methods~\cite{ZhangYu2018, Gjesteby2016}, and only more recent work has begun to consider robustness under mixed or compound degradation settings~\cite{Liang2025}. Second, although cascade architectures, dual-domain models, and diffusion-based methods have all become increasingly prominent~\cite{Schlemper2018, Wang2020, Zhang2022, Song2022, Chung2023, Liu2023DOLCE, DM4CT2025}, no published study, to the best of our knowledge, isolates cascade coupling strategy (independent, end-to-end, or residual) as the sole experimental variable under a controlled compound corruption distribution. The contribution of \Cref{chap:corruption,chap:architectures,chap:experiments} is to address precisely that gap.
\begin{table}[!htb]
    \centering
    \caption[Positioning of representative CT restoration methods.]{Positioning of representative CT image-restoration methods relative to the present thesis.
    ``Compound'' denotes joint Poisson + motion + ring + metal corruption. ``Domain'' refers to the data representation in which the network operates.}
    \label{tab:related_work}
    \footnotesize
    \setlength{\tabcolsep}{3pt}
    \begin{tabular}{@{}l l l l c@{}}
        \toprule
        Method        & Corruption regime    & Architecture           & Domain           & Ref. \\
        \midrule
        BM3D          & Gaussian noise       & Non-local filter       & Image            & \cite{Dabov2007}      \\
        DnCNN         & Gaussian noise       & Plain CNN + residual   & Image            & \cite{Zhang2017}      \\
        RED-CNN       & Poisson (LDCT)       & Encoder--decoder       & Image            & \cite{Chen2017}       \\
        FBPConvNet    & Sparse-view CT       & U-Net post-FBP         & Image            & \cite{Jin2017}        \\
        WGAN-VGG      & Poisson (LDCT)       & GAN + perceptual loss  & Image            & \cite{Yang2018WGAN}   \\
        EDRAM-Net     & Poisson (LDCT)       & ED + attention         & Image            & \cite{Komolafe2024}   \\
        U-Net-MA      & Poisson (LDCT)       & U-Net + attention      & Image            & \cite{Zhang2023}      \\
        CNN-MAR       & Metal                & CNN, dedicated MAR     & Image            & \cite{ZhangYu2018}    \\
        DuDoNet       & Metal                & Dual-domain network    & Sinogram + image & \cite{Lin2019DuDoNet} \\
        SIRNet        & Sparse-view + LDCT   & Dual-domain cascade    & Sinogram + image & \cite{Wang2020}       \\
        DSS-Net       & Sparse-view (LDCT)   & Dual-domain cascade    & Sinogram + image & \cite{Zhang2022}      \\
        MoDL          & General inverse      & Unrolled model-based   & Image            & \cite{Aggarwal2019MoDL} \\
        VarNet        & MRI                  & Unrolled variational   & k-space + image  & \cite{Hammernik2018VarNet} \\
        ArtDA         & Mixed artifacts      & Domain adaptation      & Image            & \cite{Liang2025}      \\
        Score-based MRI/CT & Sparse-view (LDCT) & Score-based prior + DC & Image            & \cite{Song2022}       \\
        DPS           & General inverse      & Diffusion posterior    & Image            & \cite{Chung2023}      \\
        DOLCE         & Limited-angle CT     & Conditional diffusion  & Image            & \cite{Liu2023DOLCE}   \\
        DM4CT         & LDCT/sparse/limited  & Diffusion benchmarks   & Image            & \cite{DM4CT2025}      \\
        \midrule
        \textbf{This thesis} & \textbf{Compound} & \textbf{Cascade study} & \textbf{Image} & ---             \\
        \bottomrule
    \end{tabular}
\end{table}
